\chapter{绪论}

代数几何是利用代数研究几何对象的学科，其中主要的研究对象是\bluebf{代数簇(Algebraic Variety)}.
\begin{definition}{$\mathbb{C}^n$中的代数簇}
    在$\mathbb{C}^n=\{(z_1,z_2,\cdots,c_n)|z_i\in\mathbb{C}\}$中，有一组复系数多项式$f_i\in\mathbb{C}[z_1,z_2,\cdots,z_n]$，则把它们的共同零点集$Z=\left\{z\in\mathbb{C}^n|f_i(z)=0,\forall\;1\leqslant i \leqslant n\right\}\subseteq \mathbb{C}^n$称为\textbf{代数簇}.
\end{definition}
为了理解这一抽象的概念，我们可以考虑以下结构的代数簇.\par
\begin{instance}
    若$f_i=\sum_{i=1}^n{a_{ij}z_j}=b_i,\,i=1,2,\cdots,m$.
    它的代数簇就是这$m$个多项式方程组成的方程组的解，将其写成矩阵形式，也就是：
    \[
    Z=\{z\in\mathbb{C}^n|Az=b\}\;,\text{其中}A=\begin{pmatrix}
        a_{11} & a_{12} & \cdots & a_{1n} \\
        \vdots & \vdots & \ddots & \vdots \\
        a_{m1} & a_{m2} & \cdots & a_{mn} 
    \end{pmatrix},\,b=\begin{pmatrix}
        b_1 \\
        b_2 \\
        \vdots \\
        b_n
    \end{pmatrix}
    \]
    因此，这一情况下的代数簇就是上述矩阵方程的解空间.\par
    \begin{enumerate}[label=case\;\arabic*]
        \item 当$b=0$时，$Ax=0$，$Z=\{z\in\mathbb{C}^n| Ax=0\}\subseteq\mathbb{C}^n$是矩阵$A$的零空间.此时$Z$构成了一个维数$\operatorname{dim}Z=n-\operatorname{rk}(A)$的线性子空间.\par
        反过来，既然代数簇可以是线性子空间，那么$V\subseteq \mathbb{C}^n$的线性子空间能否实现为一个代数簇.\par
        不妨设$V$的基为$e_1,e_2,\cdots,e_n$，由基之间线性无关可知$\exists\,a_1,a_2,\cdots a_n\text{不全为零}$使得$a_1e_1+a_2e_2+\cdots+a_ne_n=0$，利用其系数可构造一个线性方程$a_1z_1+a_2z_2+\cdots+a_nz_n=0$，记其为$f_1$.\par
        由于$\mathbb{C}^n$是有限维线性空间，我们可以在其中取遍满足$\sum_{i=1}^n{a_ie_i}$且线性无关的$(a_1,a_2,\cdots,a_n)$，并用它们构造出若干个线性方程$f_1,f_2,\cdots,f_m$.\par
        最终我们可以证明$V=Z=\{z\in\mathbb{C}^n|f_i(z)=0,\;1\leqslant i\leqslant m\}$，可以实现为一个代数簇.
        \item 当$b\neq 0$时，$Az=b$有解$\Leftrightarrow \operatorname{rk}(A;b)=\operatorname{rk}(b)$.\par
        如果$Az=b$有解$z_0$，即特解，则通解可通过$A(z-z_0) =0$求得.\par
    \end{enumerate}
\end{instance}

\begin{instance}
    考虑二次曲线$\{(x,y)\in\mathbb{C}^2|x^2+y^2=1\}\subseteq\mathbb{C}^2$，它实际就是$x^2+y^2=1$对应的代数簇.在平面坐标系上，它是单位圆，在$\mathbb{C}^2$上，我们无法直观画出其图像，作为替代，需要找到其参数化曲线.\par
    \begin{wrapfigure}{l}{0.3\textwidth}
        \centering
        \vspace{-1em}
        \includegraphics[width=0.3\textwidth]{fig/1.1.png}
        \vspace{-6em}
    \end{wrapfigure}
    先取这个曲线上的一点$(0,1)$，除了该点外，圆上所有其他的点都可以视作是过$(0,1)$的直线与圆的另一个交点.对于一个过定点的动直线，只需要一个参数就可以描述它.在这里，我们选择使用斜率来刻画这条直线(当然这样会漏掉$(0,-1)$这种情形).\par
    如果设这条直线的斜率为$\lambda$，那么直线方程就可以写成$y=\lambda x+1$，将其与单位圆的方程联立就可以得到其交点坐标为$(x,y)=\left(-\frac{2\lambda}{\lambda^2+1},\frac{1-\lambda^2}{\lambda^2+1}\right)$.\par
    现在稍微变换一下方程，考虑方程$x^2+y^2=-1$和$x^2+y^2=0$，前者在$\mathbb{R}^2$上无解，后者有唯一解$(0,0)$.\par
    但如果在$\mathbb{C}^2$上，可以对后者作因式分解得到$x^2+y^2=x^2-(-y^2)=(x+\sqrt{-1}y)(x-\sqrt{-1}y)=0$，它在$\mathbb{C}^2$上是两条直线的并.\par
    如果我们进一步拓展问题，考虑满足$X^2+Y^2=Z^2$的点，由于上述方程是齐次的，所以若$(X,Y,Z)$满足上式，则$(\mu X,\mu Y,\mu Z)$也满足.可以使用初等数论的工具得到该方程的通解的形式.\par
    下一个问题是方程$2X^2+Y^2=5Z^2$是否有非负整数解？\par
    首先可以观察到，如果$5|X,5|Y \Rightarrow 25|5Z^2 \Rightarrow 5|Z^2 \Rightarrow 5|Z$.因此我们可以不断使用$\left(\frac{X}{5},\frac{Y}{5},\frac{Z}{5}\right)$代替$(X,Y,Z)$，直到某一步后$X,Y$中至少有一个不能被$5$整除.\par
    在上面的化简步骤完成后，我们就可以在$\mathbb{Z}/5\mathbb{Z}$上看这个问题，在$\mathbb{Z}_5$上，方程转化为$\bar{2}\overline{X}^2 + \overline{Y}^2=\bar{0}$.
    由于$\mathbb{Z}_5$是一个有限域，因此可以写出$\overline{X},\overline{Y}$的所有可能的取值，即$\overline{X},\overline{Y}\in\{\bar{0},\bar{1},\bar{4}\}$.\par
    若$\overline{Y}=\bar{0} \Rightarrow \overline{X}\neq \bar{0}$，否则就会和预设的$X,Y$至少其一不被5整除矛盾，但当$\overline{X}\neq 0$时，一一验证会发现不存在满足条件的$\overline{X}$.\par
    对$\overline{Y}=\bar{1},\bar{4}$也同样作验证，结果是始终不存在解，因此该方程没有非零整数解.在这一过程中，我们使用到了一个重要的方法：\textbf{reduction mod p}.
\end{instance}

在绪论的最后，复习二次型相关的知识.\par

\begin{definition}{二次型}
    设$Ｖ$是一个向量空间，$K$是一个数域，则
   \begin{align*}
       Q:V\times V &\to K \\
       (u,v) &\mapsto Q(u,v) \\
       Q(u,v)&=Q(v,u)
   \end{align*}称为$V$上的\textbf{二次型}.
\end{definition}
在$\mathbb{C}^n$上，可以使用矩阵在表示一个二次型.对于每一个二次型对应的函数，我们总可以找到一个对称矩阵
$A=\begin{pmatrix}
    a_{11} & a_{12} & \cdots & a_{1n} \\
    \vdots & \vdots & \ddots & \vdots \\
    a_{n1} & a_{n2} & \cdots & a_{nn} \\
\end{pmatrix}$，使得$f(x)=xAx^T$，因此$f(x)$是一个对称多项式.\par
下一个问题是如何分类不同的二次型，由于一个二次型可以用对称矩阵表示，因此分类二次型实际上就是在分类对称矩阵.\par
我们可以对对称矩阵作合同变换，即如果$A,B$是两个对称矩阵，并且存在可逆矩阵$P$，使得$B=PAP^T$，就称$A$和$B$是合同的，在不同数域下，对称矩阵会合同于不同的矩阵.\par
在$\mathbb{C}^n$上，可以证明任意对称矩阵$A$一定会合同于
$\begin{pmatrix}
    I_{i} & \\
    & O_n
\end{pmatrix}$，单位矩阵的阶数等于$\mathrm{rk}(A)$，此时二次型为$f=x_1^2+x_2^2+\cdots+x_i^2$.\par
在$\mathbb{R}^n$上，可以证明任意对称矩阵$A$一定会合同于$\begin{pmatrix}
    I_a & & \\
    & -I_b & \\
    & & O_n \\
\end{pmatrix}$，此时二次型为$f=x_1^2+x_2^2+\cdots+x_a^2-y_1^2-y_2^2-\cdots-y_b^2$，$a$和$b$被分别称作\bluebf{正惯性指数}和\bluebf{负惯性指数}.
